% Numeracao de eq./fig./tab. por secao do apendice (do modelo)
\renewcommand\theequation{\Alph{section}\arabic{equation}}
\counterwithin*{equation}{section}
\renewcommand\thefigure{\Alph{section}\arabic{figure}}
\counterwithin*{figure}{section}
\renewcommand\thetable{\Alph{section}\arabic{table}}
\counterwithin*{table}{section}

\begin{appendices}

% Formatacao dos apendices (do modelo)
\renewcommand{\thesection}{Ap\^endice \Alph{section}}
\renewcommand{\thesubsection}{\Alph{section}.\arabic{subsection}}

%--- Apêndice A: procedimento correspondente aos notebooks atuais ---%
\section{Verificação da margem de durabilidade em uma idade}
\label{app:exemplo_passo}

Este apêndice explicita a sequência de cálculo de uma resposta do Exemplo~2. Considera-se o ponto nominal $f_c=25$~MPa, $UR=55\%$ e $c_{\mathrm{nom}}=25$~mm, com início da exposição em 1980, cimento CP~II~F, exposição externa protegida, $ad=0\%$ e cenário SSP2-4.5. Para uma realização latente $i$, calculam-se
\begin{equation}
    f_{c,i}^{\ast}=25\,\omega_{f_c,i},\qquad
    UR_i^{\ast}=0{,}55\,\omega_{UR,i},\qquad
    c_i=25\,\omega_{c,i}.
\end{equation}
A umidade é expressa como fração no modelo, enquanto $f_c$ e $c$ permanecem em MPa e mm. Para cada idade da malha auxiliar de dez anos, obtém-se a concentração no ano-calendário correspondente e avalia-se a Equação~\eqref{eq:y}. O máximo acumulado dessas profundidades define $\tilde{y}_i(t)$ nos pontos da malha.

Para uma consulta aos 25 anos, correspondente a 2005, o simulador interpola entre as profundidades acumuladas aos 20 e 30 anos:
\begin{equation}
    y_i(25)=\tilde{y}_i(20)+\frac{25-20}{30-20}
    \left[\tilde{y}_i(30)-\tilde{y}_i(20)\right],\qquad
    g_i(25)=c_i-y_i(25).
\end{equation}
O resultado é uma margem em milímetros. Um valor $g_i(25)\leq0$ identifica que a frente alcançou a armadura. Para a política principal, $g_i(25)\leq5$~mm identifica que o critério preventivo já foi atingido nessa idade, dada a monotonicidade do simulador.

Repetindo esse procedimento para as $2.500$ réplicas de um mesmo ponto, obtém-se a amostra condicional de $g$ à qual a GLD é ajustada. A repetição para os pontos de treinamento e de validação produz os conjuntos usados na construção das PCEs. Este cálculo pertence à geração dos dados de referência; a etapa posterior de vida remanescente utiliza as distribuições previstas pelas redes.

% Material anterior de flexão preservado apenas como rascunho; não integra
% o simulador de despassivação nem o texto compilado deste artigo.
\iffalse
%--- Apendice A ---%
\section{Exemplo de verifica\c{c}\~ao para um passo de tempo}
\label{app:exemplo_passo}

% Material DIDATICO preservado do grupo (antiga Section_5). Detalhado
% demais para o corpo de um artigo Q1, mas excelente para a disserta\c{c}\~ao.
% Mantido em portugues e com os numeros originais dos alunos.

Este apêndice detalha, para um único passo de tempo e uma única simulação, o procedimento de cálculo da margem de segurança $g(t)$ descrito na Seção~\ref{subsec:resistencia_residual}, a fim de tornar transparente a sequência de operações do simulador. A viga considerada segue os parâmetros da Seção~\ref{subsec:exemplo2}, com cobrimento nominal $c_{\mathrm{nom}} = 25$~mm, diâmetro nominal das barras $d_{s0} = 12{,}5$~mm e coeficiente de carbonatação $k = 6$~mm/$\sqrt{\text{ano}}$.

Considere o passo de tempo $t = 40$ anos, com temperatura ambiente $T = 23{,}56\,^\circ$C, taxa de corrosão a 20\,$^\circ$C igual a $i_{\mathrm{corr},20} = 0{,}431\,\mu$A/cm$^2$ e início da corrosão em $t_{\mathrm{ini}} = 25$ anos.

\subsection{Verifica\c{c}\~ao da carbonata\c{c}\~ao}

\begin{equation}
y_{\mathrm{carb}}(40) = 6\sqrt{40} = 37{,}95~\text{mm}.
\end{equation}

Como $y_{\mathrm{carb}}(40) > c_{\mathrm{nom}} = 25$~mm, considera-se que já ocorreu o início da corrosão.

\subsection{Taxa de corros\~ao ajustada}

Pela Equação~\eqref{eq:i_cor_t}, com $T > 20\,^\circ$C ($K = 0{,}073$):

\begin{equation}
i_{\mathrm{corr}}(23{,}56) = 0{,}431\,\left[\,1 + 0{,}073\,(23{,}56 - 20)\,\right] = 0{,}542~\mu\text{A/cm}^2.
\end{equation}

\subsection{Perda de di\^ametro e \'area efetiva}

Com $t_{\mathrm{corr}} = 40 - 25 = 15$ anos, a Equação~\eqref{eq:d_cor} fornece

\begin{equation}
d_{s}(40) = 12{,}5 - 0{,}0232 \times 0{,}542 \times 15 = 12{,}3114~\text{mm},
\end{equation}

de modo que a área da barra corroída, pela Equação~\eqref{eq:as_t} (por barra), é

\begin{equation}
\frac{\pi}{4}\,d_s(40)^2 = \frac{\pi}{4}\,(12{,}3114)^2 = 119{,}06~\text{mm}^2,
\end{equation}

ante $A_{s,0} = 122{,}72$~mm$^2$ da barra íntegra ($d_{s0} = 12{,}5$~mm), uma redução de $3{,}66$~mm$^2$ por barra.

\subsection{Momento resistente degradado e coeficiente de ader\^encia}

Recompondo o equilíbrio seccional (Equação~\ref{eq:equilibrio}) com a área corroída, obtém-se $M_{Rd}(40) = 34{,}82$~kN$\cdot$m. O coeficiente redutor de aderência, pela Equação~\eqref{eq:c_f}, é

\begin{equation}
C_f = \frac{5{,}0}{\,12{,}5^{0{,}54}\left[(0{,}542/1000)\cdot 15 \cdot 365\right]^{0{,}19}} = 0{,}873,
\end{equation}

resultando, pela Equação~\eqref{eq:m_res}, em

\begin{equation}
M_{Rd,\mathrm{cor}}(40) = C_f \cdot M_{Rd}(40) = 0{,}873 \times 34{,}82 = 30{,}39~\text{kN}\cdot\text{m}.
\end{equation}

\subsection{Margem de seguran\c{c}a}

Com os momentos solicitantes $M_{Gk} = 12{,}0$~kN$\cdot$m e $M_{Qk} = 8{,}0$~kN$\cdot$m, o momento solicitante total é $M_S = 20{,}0$~kN$\cdot$m, e a margem de segurança (Equação~\ref{eq:funcao_g}) torna-se

\begin{equation}
g(40) = M_{Rd,\mathrm{cor}}(40) - M_S = 30{,}39 - 20{,}0 = 10{,}39~\text{kN}\cdot\text{m}.
\end{equation}

O resultado positivo indica que, aos 40 anos, mesmo considerando a corrosão, a seção ainda possui capacidade resistente superior às solicitações atuantes. Este cálculo é repetido para todos os instantes da simulação e para todas as amostras estocásticas, constituindo o simulador cujo custo o emulador da Seção~\ref{sec:emuladores} busca contornar.

\fi

%--- Apendice B (opcional) ---%
\section{Material suplementar}
\label{app:suplementar}

\falta{Incorporar os diagnósticos complementares das cinco idades de treinamento, a identificação dos pontos usados nas comparações distribucionais, os descartes da base temporal e os resultados de verificação da resolução e da monotonicidade das trajetórias.}

\end{appendices}
